(a) Every constructible subset is obtained by finitely many unions, intersections, and complements from finitely many open subsets $U_1,\ldots,U_n$. The sets obtained by choosing either $U_i$ or $X\setminus U_i$ for every $i$ and intersecting the choices are pairwise disjoint and locally closed. Their unions include every Boolean combination of the $U_i$, proving the required decomposition. Conversely, every locally closed subset is constructible, as is any finite union of them.

(b) Write the constructible subset as a finite union of locally closed subsets. If it is dense, irreducibility implies that one of these subsets is dense. Writing that subset as $U\cap F$ with $F$ closed, density forces $F=X$, so it is a nonempty open and contains the generic point. Conversely, a subset containing the generic point is dense.

(c) Closed subsets are constructible and stable under specialization. Conversely, let $S$ have these two properties and put $F=\overline S$. For each irreducible component $F_i$ of $F$, the subset $S\cap(F_i\setminus\bigcup_{j\ne i}F_j)$ is dense in that nonempty open subset of $F_i$, since $S$ is dense in $F$. It is constructible, so (b) implies that it contains the generic point of $F_i$. Specialization stability then gives $F_i\subseteq S$ for every $i$. Hence $S=F$ is closed. Taking complements gives the corresponding assertion for open subsets.

(d) Inverse images preserve finite intersections and complements, and the inverse image of an open is open. Thus they preserve constructible subsets.
